\chapter{Baire's Theorem}\label{ch:baire}

\begin{center}\itshape
Pick a continuous function without looking. Does it have a tangent anywhere?
Pick a real number the same way. Do fractions come very close, or do they keep a gap?
Pick a string of digits. Are the digits normal?
Measure is one reading of the word ``typical''. This chapter gives a second reading. The two readings often split.
\end{center}

On 18 July 1872, Karl Weierstrass read a note to the Berlin Academy. The function in that note is continuous at every real $x$, and the difference quotient at $x$ refuses to settle. The graph keeps a wiggle at every scale, so it has no tangent line.

\portrait{weierstrass}{Karl Weierstrass (1815--1897). Portrait by Conrad Fehr, 1895. Public domain, via Wikimedia Commons.}

Charles Hermite wrote to Stieltjes on 20 May 1893: ``I turn away with fright and horror from this lamentable plague of continuous functions which have no derivatives.'' Henri Poincar\'e took the same examples up in \emph{Science et m\'ethode} (1908). He wrote that logic sometimes produces monsters, and that functions were being invented on purpose to defeat the reasoning of our fathers.

In 1931 Stefan Banach and Stefan Mazurkiewicz stopped building examples one by one. They took the whole space of continuous functions on $[0,1]$ and showed that the functions with a derivative at even a single point form a meagre set, in the sense defined below. Weierstrass's example is what a continuous function usually is. The functions with a derivative, the ones a calculus course draws, form the thin set.

The tool is older than that paper. Ren\'e-Louis Baire's thesis of 1899 proves that a complete metric space cannot be written as a countable union of nowhere dense sets. The word ``most'' then has a meaning, and the meaning does not mention lengths.

\portrait{baire}{Ren\'e-Louis Baire (1874--1932). Photograph, about 1900, author unknown. Public domain, via Wikimedia Commons.}

The proof shows that a point with the property exists because almost every point has it. You finish the argument without writing down a single example. A haystack that is mostly needles is the same kind of surprise: you know the needles are there before you have held one. Three theorems in functional analysis sit on this fact. They are the uniform boundedness principle, the open mapping theorem, and the closed graph theorem, all in Volume~8.
% MATHFIX: those three theorems are named as later uses in Volume~8. They are not proved in this chapter, and no proof below treats them as premises.
Inside this chapter the same fact does three jobs. A typical continuous function has a derivative nowhere. The Liouville numbers form a set of measure zero that is still typical in Baire's sense. No function from $\R$ to $\R$ is continuous at the rationals and discontinuous at the irrationals.

\section{Nowhere dense and meagre sets}

Here is the picture on the line, before the names. The integers $\Z$ are closed in $\R$, and no open interval sits inside $\Z$. The rationals $\Q$ behave differently: every open interval contains rationals, and the closure of $\Q$ is the whole line. The definition splits those two pictures.

\begin{definition}
Let $A$ be a subset of a topological space $X$. Call $A$ \emph{nowhere dense} when the interior of its closure is empty. Call $A$ \emph{meagre}, or \emph{of first category}, when $A$ sits inside a countable union of nowhere dense sets. The complement of a meagre set is \emph{comeagre}. A set is a $G_\delta$ when it equals a countable intersection of open sets, and an $F_\sigma$ when it equals a countable union of closed sets.
\end{definition}

Nowhere dense has a useful open-set form. The closure $\bar A$ contains no nonempty open set, so if $U$ is nonempty and open, the difference $V=U\setminus\bar A$ is open and nonempty, and $V$ misses $A$. Conversely, if every nonempty open set contains a nonempty open set that misses $A$, then $\bar A$ cannot contain a nonempty open set. A finite subset of $\R$ is nowhere dense. So is $\Z$. So is the Cantor set: it is closed, and it contains no interval. Chapter~\ref{ch:compact} records that the stage $C_n$ has total length $\bigl(\frac23\bigr)^n$. Chapter~\ref{ch:measures} is where that shrinking length becomes a measure. The set $\Q$ is meagre, being a countable union of single points, and $\Q$ is dense. A subset of a meagre set is meagre: it sits in the same countable union. A countable union of meagre sets is meagre: a countable union of countable unions is still a countable union.

\begin{figure}[ht]
\centering
\fig{ch05-rw-fatcantor}
\caption{Two stages of a fat Cantor set in $[0,1]$. The charcoal bars are what remains, and the coral segments are open gaps. Stage $1$ deletes an open middle of length $\frac14$. Stage $2$ deletes an open middle of length $\frac1{16}$ from each surviving bar. Every bar that you can still see contains an interval, so the stage itself is not nowhere dense. The later stages cut a gap from every one of those bars. The intersection of all the stages is closed and contains no interval, which is what nowhere dense asks for on the line.}
\label{fig:rwfatcantor}
\end{figure}

\begin{theorem}[Baire, 1899]\label{thm:baire}
Let $X$ be a complete metric space, or a locally compact Hausdorff space, so that every point has a neighbourhood whose closure is compact. Any countable intersection of dense open subsets of $X$ is dense. Equivalently, every meagre subset of $X$ has empty interior. In particular, a nonempty open subset of $X$ is not meagre, and $X$ itself is not meagre when $X$ is nonempty.
\end{theorem}
\begin{proof}
\emph{The two wordings agree.} A set $U$ is open and dense exactly when $X\setminus U$ is closed and has empty interior, which is the closed case of nowhere dense. Let $A$ be meagre, so $A$ sits in a countable union of closed nowhere dense sets $\overline{A_n}$. In any space, $A$ has empty interior exactly when $X\setminus A$ is dense. The complement $X\setminus A$ contains
\[
\bigcap_n\bigl(X\setminus\overline{A_n}\bigr),
\]
and each factor $X\setminus\overline{A_n}$ is open and dense. If every such intersection is dense, then $X\setminus A$ contains a dense set, so $X\setminus A$ is dense, and the interior of $A$ is empty. In the other direction, take dense open sets $U_n$. The union of the closed nowhere dense sets $X\setminus U_n$ is meagre. If every meagre set has empty interior, that union has empty interior, so $\bigcap_n U_n$ is dense.

\emph{Complete metric spaces.} Let $U_1,U_2,\dots$ be open and dense, and let $G$ be open and nonempty. The goal is a point of $G\cap\bigcap_n U_n$. Density of $U_1$ makes $G\cap U_1$ open and nonempty, so that intersection contains a closed ball $\bar B(x_1,r_1)$ with $0<r_1<1$. Density of $U_2$ makes $B(x_1,r_1)\cap U_2$ contain a closed ball $\bar B(x_2,r_2)$ with $0<r_2<\frac12$. Continue. You obtain closed balls satisfying
\[
\bar B(x_{n+1},r_{n+1})\subseteq B(x_n,r_n)\cap U_{n+1},\qquad 0<r_n<\frac1n.
\]
The radii are a concrete sequence of room: less than $1$, then less than $\frac12$, then less than $\frac13$. If $m\ge n$, the nesting puts $x_m$ in $B(x_n,r_n)$, so $d(x_m,x_n)<r_n<\frac1n$. The sequence $(x_n)$ is Cauchy. Completeness gives a limit $x$. For $m\ge n$ the point $x_m$ lies in the closed ball $\bar B(x_n,r_n)$, and a closed set contains its limits, so $x$ lies in that ball. The first ball sits in $G\cap U_1$. For $n\ge2$ the construction puts $\bar B(x_n,r_n)$ inside $U_n$, and the nesting $\bar B(x_n,r_n)\subseteq B(x_1,r_1)$ keeps the ball inside $G$. Therefore $x\in G\cap U_n$ for every $n$.

\emph{Locally compact Hausdorff spaces.} First, an open neighbourhood can be shrunk until its closure is compact and still inside the open set you started with. Let $x\in U$ with $U$ open, and choose an open set $W\ni x$ whose closure $K=\bar W$ is compact. The set $K$ is compact Hausdorff, so it is normal by Proposition~\ref{prop:compactbasic}(e). Inside $K$, the point $x$ and the closed set $K\setminus(U\cap W)$ can be put in disjoint sets $A\ni x$ and $B\supseteq K\setminus(U\cap W)$, both open in $K$. Since $A$ is open in $K$, there is an open set $O$ in $X$ with $A=O\cap K$. The inclusion $W\subseteq K$ then gives $A\cap W=O\cap W$, which is open in $X$. Set $V=A\cap U\cap W$. Then $V$ is open in $X$. Also $V\subseteq W$, so $\bar V\subseteq K$. The same $V$ sits in $A$, so its closure in $K$ misses $B$, and
\[
\bar V\subseteq K\setminus B\subseteq U\cap W.
\]
A closed subset of the compact set $K$ is compact, so $\bar V$ is compact. Now copy the metric construction with these neighbourhoods in place of balls. Choose a nonempty open $V_1$ with $\overline{V_1}$ compact and $\overline{V_1}\subseteq G\cap U_1$. Then choose a nonempty open $V_2$ with $\overline{V_2}\subseteq V_1\cap U_2$, and continue. The compact sets $\overline{V_n}$ are nested and nonempty. They are closed subsets of the compact set $\overline{V_1}$, so the finite-intersection property supplies a point common to all of them. That point lies in $G\cap\bigcap_n U_n$.
\end{proof}

Each ball, or each neighbourhood $V_n$, is chosen after the previous one. The proof uses countably many choices of that dependent kind. It does not use the full axiom of choice.

Both hypotheses are needed, and $\Q$ shows why. With the usual distance, $\Q$ is a countable union of singletons, and each singleton is nowhere dense in $\Q$, so $\Q$ is meagre in itself. The space $\Q$ is neither complete nor locally compact. A meagre set can still be large in another sense. It can be dense, as $\Q$ is. It can be uncountable, as the Cantor set is. Section~\ref{sec:categorymeasure} will show that it can even have full measure.

\begin{figure}[ht]
\centering
\fig{ch05-rw-lines}
\caption{Five horizontal lines in a window of the plane, drawn in coral at the heights $0,\frac12,1,\frac32,2$. Each line is closed, and no disc is contained in a line, so each line is nowhere dense. The same is true of the line at every rational height. There are countably many rationals, so the union of those lines is meagre in $\R^2$. The charcoal region is the complement. The disc centred on the line $y=1$ still contains charcoal points, and the small charcoal disc between two lines is an open piece of the complement. Theorem~\ref{thm:baire} applies because $\R^2$ is complete, so this complement is comeagre: it is a residual set.}
\label{fig:rwlines}
\end{figure}

\begin{corollary}\label{cor:bairecor}
\begin{enumerate}[label=(\alph*)]
\item If a complete metric space is a countable union of closed sets $F_n$, then some $F_n$ has nonempty interior.
\item A dense $G_\delta$ in a complete metric space is comeagre. A countable intersection of dense $G_\delta$ sets is again a dense $G_\delta$.
\item The line $\R$ is uncountable, and so is every nonempty interval.
\end{enumerate}
\end{corollary}
\begin{proof}
(a) If every $F_n$ had empty interior, then $X$ would be a countable union of closed nowhere dense sets, hence meagre in itself. Theorem~\ref{thm:baire} forbids that when $X$ is nonempty.

(b) A dense $G_\delta$ has the form $\bigcap_n U_n$ with each $U_n$ open. Each $U_n$ contains the dense $G_\delta$, so each $U_n$ is dense. The complement is the union of the closed nowhere dense sets $X\setminus U_n$, so the $G_\delta$ is comeagre. A countable intersection of countable intersections of dense open sets is one intersection of dense open sets. Theorem~\ref{thm:baire} makes that intersection dense, and it is a $G_\delta$ by construction.

(c) A singleton in $\R$ is closed and has empty interior. Let $I$ be a nonempty interval. If $I$ is closed it is complete. If it is not closed, any closed subinterval with positive length is complete. If $I$ were countable, that complete piece would be a countable union of singletons, contradicting (a).
\end{proof}

Part~(c) is Cantor's theorem again, Theorem~\ref{thm:cantor}, reached from completeness instead of from the diagonal. Exercise~\ref{ex:perfect} carries the same idea past the line.

\begin{definition}
Let $X$ be a complete metric space. A property of points of $X$ is \emph{generic} when the set of points with the property is comeagre, so the property holds at least on a dense $G_\delta$. A point of that comeagre set is called \emph{typical}.
\end{definition}

Corollary~\ref{cor:bairecor}(b) lets you stack generic properties. Countably many comeagre sets have a comeagre intersection, so a typical point has all of those properties at once. The same stacking is familiar from probability, where countably many almost sure events happen together.

\section{Nowhere differentiable functions}

\subsection*{An explicit example}

One monster can be written as a series. Let $\phi(x)=\operatorname{dist}(x,\Z)$. This is the $1$-periodic tent map: on $[0,1]$ it rises with slope $1$ from $0$ to $\frac12$ and falls with slope $-1$ from $\frac12$ to $1$. Takagi's function, from 1903, is the uniform series
\[T(x)=\sum_{n\ge0}\frac{\phi(2^nx)}{2^n}.\]
The general term is at most $2^{-n-1}$, because $\phi\le\frac12$. The comparison $\sum 2^{-n-1}<\infty$ gives uniform convergence, and a uniform limit of continuous functions is continuous by Proposition~\ref{prop:complete}(c). So $T$ is continuous on $\R$.

\begin{figure}[ht]
\centering
\fig{ch05-takagi}
\caption{Takagi's function $T$ on $[0,1]$, and two straightened zooms. The lower pictures show the graph over $[\frac14,\frac38]$ and over $[\frac5{16},\frac{21}{64}]$, of lengths $\frac18$ and $\frac1{64}$, after the chord has been subtracted and both axes have been scaled to the same frame. On a dyadic interval of length $2^{-k}$ the identity $T(u+2^{-k}t)=T(u)+2^{-k}T(t)+(\text{affine in }t)$ says that the straightened zoom is a copy of the whole graph. Differentiability at a point would make some zoom look like a straight line. None of these zooms does.}
\label{fig:takagi}
\end{figure}

The next lemma is the test the series will fail. If a derivative exists, difference quotients over shrinking brackets have to approach it, even when the point sits strictly between the ends.

\begin{lemma}\label{lem:squeeze}
Suppose $f$ is differentiable at $x$. Let $u_k\le x\le v_k$ with $u_k<v_k$ and $v_k-u_k\to0$. Then
\[\frac{f(v_k)-f(u_k)}{v_k-u_k}\to f'(x).\]
\end{lemma}
\begin{proof}
First take $u_k<x<v_k$, and set $\lambda=\frac{v_k-x}{v_k-u_k}$. Then $\lambda\in\,]0,1[$, and
\[\frac{f(v_k)-f(u_k)}{v_k-u_k}=\lambda\,\frac{f(v_k)-f(x)}{v_k-x}+(1-\lambda)\,\frac{f(x)-f(u_k)}{x-u_k}.\]
The right-hand side is a convex combination of the two difference quotients of $f$ at $x$. Differentiability puts both of those quotients within $\varepsilon$ of $f'(x)$ once $k$ is large, and a convex combination of two such numbers is still within $\varepsilon$ of $f'(x)$. If instead $x=u_k$ or $x=v_k$, the quotient on the left is already a difference quotient at $x$, so it tends to $f'(x)$ directly.
\end{proof}

\begin{theorem}\label{thm:takagi}
Takagi's function $T$ is continuous on $\R$, and $T$ is differentiable at no point of $\R$.
\end{theorem}
\begin{proof}
Continuity was checked above. The function $\phi$ has period $1$, and so does every term $\phi(2^nx)$, so $T$ has period $1$. It is enough to rule out a derivative at each $x\in[0,1[$. Fix $k\ge1$, and let $u_k=j2^{-k}\le x<v_k=(j+1)2^{-k}$ be the dyadic interval of length $2^{-k}$ that contains $x$.

Look at one term of the series on that interval. If $n\ge k$, then $2^nu_k$ and $2^nv_k$ are integers, so $\phi$ vanishes at both ends and the $n$-th term contributes $0$ to the difference $T(v_k)-T(u_k)$. If $n<k$, the map $x\mapsto 2^{-n}\phi(2^nx)$ is affine of slope $\pm1$ on each interval $[i2^{-n-1},(i+1)2^{-n-1}]$. The dyadic interval $[u_k,v_k]$ has length $2^{-k}$ and $k\ge n+1$, so it sits inside one of those pieces. The difference quotient of that term is therefore some sign $\varepsilon_n\in\{-1,+1\}$. Adding the finitely many terms with $n<k$ gives
\[s_k=\frac{T(v_k)-T(u_k)}{v_k-u_k}=\sum_{n<k}\varepsilon_n,\qquad\varepsilon_n\in\{-1,+1\}.\]
The next quotient drops or adds exactly one sign: $s_{k+1}-s_k=\varepsilon_k=\pm1$. A sequence that jumps by $1$ at every step cannot converge. Lemma~\ref{lem:squeeze} then says that $T'(x)$ does not exist.
\end{proof}

Weierstrass's series was a sum of cosines,
\[W(x)=\sum_na^n\cos(b^n\pi x),\]
where $0<a<1$, the frequency $b$ is an odd integer, and $ab>1+\frac{3\pi}2$. Hardy, in 1916, showed that the weaker condition $ab\ge1$ already forces $W$ to be nowhere differentiable. That improvement is quoted from Hardy. This chapter does not prove it.
% MATHFIX: Hardy's improvement \(ab\ge 1\) is quoted from Hardy (1916) and is not proved here.
Cosines have no corners, which is why his argument is longer than the sign-chasing above. Figure~\ref{fig:weierstrass} uses $a=\frac12$ and $b=3$, so $ab=\frac32$. This pair meets Hardy's bound $ab\ge1$. It misses Weierstrass's bound, since $1+\frac{3\pi}2$ is larger than $\frac32$.

\begin{figure}[ht]
\centering
\fig{ch05-weierstrass}
\caption{Partial sums of $\sum_{n\ge0}a^n\cos(b^n\pi x)$ on $[0,1]$, with $a=\frac12$ and $b=3$. The light graph uses two terms, the middle graph four, and the dark graph six. Term $n$ oscillates $b^n=3^n$ times as fast as the first cosine, and its height is $a^n=2^{-n}$. Differentiating brings a factor $b^n\pi$, so the slopes of the new term are about $\frac32$ times as large as those of the previous term. The wiggles get steeper as they get thinner.}
\label{fig:weierstrass}
\end{figure}

\subsection*{Most continuous functions}

Takagi's series is one explicit curve. Banach and Mazurkiewicz did not write a series. They measured, in Baire's sense, the set of functions in $C([0,1])$ that behave at even one point. That set is meagre. If you sample a continuous function in the sense of category, you do not see a derivative.

\begin{theorem}[Banach, Mazurkiewicz, 1931]\label{thm:banachmaz}
Equip $C([0,1])$ with the supremum norm. The set of functions that are differentiable at some point of $[0,1]$, allowing a one-sided derivative at an endpoint, is meagre. A typical function in $C([0,1])$ is therefore differentiable at no point, and the nowhere differentiable functions are dense in $C([0,1])$.
\end{theorem}
\begin{proof}
For each integer $n\ge2$, let $E_n$ be the set of $f\in C([0,1])$ for which some point $x\in[0,1-\frac1n]$ satisfies
\[|f(x+h)-f(x)|\le nh\quad\text{for all }h\in\,]0,\tfrac1n].\]
The point $x$ is a witness that the difference quotients to the right of $x$, at least out to distance $\frac1n$, stay bounded by $n$.

\emph{A finite right derivative puts $f$ in some $E_n$.} Suppose the right derivative of $f$ exists and is finite at some $x\in[0,1[$. The definition of that derivative gives $C$ and $\delta\in\,]0,1-x]$ such that $|f(x+h)-f(x)|\le Ch$ whenever $0<h\le\delta$. On the range $\delta\le h\le1-x$ the crude bound $|f(x+h)-f(x)|\le2\|f\|_\infty$ is enough, and $h\ge\delta$ turns it into
\[|f(x+h)-f(x)|\le\frac{2\|f\|_\infty}{\delta}\,h.\]
Choose an integer $n\ge\max\bigl(C,\frac{2\|f\|_\infty}{\delta},\frac1{1-x}\bigr)$. Then $x\le1-\frac1n$, and every $h\in\,]0,\frac1n]$ falls under one of the two bounds, so $x$ is a witness for $E_n$.

\emph{The set $E_n$ is closed.} Let $f_k\in E_n$ converge uniformly to $f$, and let $x_k$ be a witness for $f_k$. The witnesses lie in the compact interval $[0,1-\frac1n]$, so a subsequence, still written $x_k$, converges to some $x$ in that interval. Fix $h\in\,]0,\frac1n]$. Then $|f_k(x_k+h)-f_k(x_k)|\le nh$. Uniform convergence and continuity of $f$ pass to the limit along the subsequence. Explicitly,
\[
|f_k(x_k)-f(x)|\le\|f_k-f\|_\infty+|f(x_k)-f(x)|,
\]
and the same estimate with $x_k+h$ in place of $x_k$ gives $f_k(x_k+h)\to f(x+h)$. Therefore $|f(x+h)-f(x)|\le nh$. The limit $x$ is a witness for $f$, and $f\in E_n$.

\emph{The interior of $E_n$ is empty.} Fix $f\in C([0,1])$ and $\varepsilon>0$. Corollary~\ref{cor:weierstrass} supplies a polynomial $p$ with $\|f-p\|_\infty<\frac\varepsilon2$. The derivative $p'$ is continuous on the compact interval $[0,1]$, so it is bounded, say $|p'|\le M$. Theorem~\ref{thm:mvt} turns that bound into the Lipschitz estimate $|p(x+h)-p(x)|\le Mh$. Build a sawtooth $s$: continuous, piecewise affine, with values in $[0,\frac\varepsilon2]$, and with slopes $\pm L$ on pieces shorter than $\frac1n$, where the slope size satisfies $L>n+M$. From any $x\in[0,1-\frac1n]$ you can then step a small $h\in\,]0,\frac1n]$ and stay inside the piece immediately to the right of $x$. Set $g=p+s$. The triangle inequality gives $\|g-f\|_\infty<\varepsilon$. On that short step,
\[
\begin{aligned}
|g(x+h)-g(x)|&\ge|s(x+h)-s(x)|-|p(x+h)-p(x)|\eqbr
&\ge(L-M)h>nh.
\end{aligned}
\]
No point $x$ is a witness for $g$, so $g\notin E_n$. Every ball around $f$ meets the complement of $E_n$.

Each $E_n$ is closed and has empty interior, so each $E_n$ is nowhere dense. Every $f$ with a finite right derivative at some point of $[0,1[$ lies in $\bigcup_n E_n$, a meagre set. The reflection $f\mapsto f(1-\,\cdot\,)$ is an isometry of $C([0,1])$, and it swaps left derivatives with right derivatives. The functions with a finite left derivative somewhere in $]0,1]$ are therefore meagre as well. A function that is differentiable at an interior point has both one-sided derivatives there, and a function that is differentiable at an endpoint has the one-sided derivative the endpoint allows. Either way, the function lies in one of the two meagre sets. The space $C([0,1])$ is complete, so Theorem~\ref{thm:baire} says that the complement is comeagre. A comeagre subset of a complete metric space is dense, and the complement is exactly the set of nowhere differentiable functions.
\end{proof}

\begin{beauty}
The argument never exhibits a formula for a single wild function. It shows something stronger about location. Every ball in $C([0,1])$ contains nowhere differentiable functions, because a comeagre set is dense. Polynomials are dense as well, by Corollary~\ref{cor:weierstrass}, and every polynomial is differentiable, so the meagre set in Theorem~\ref{thm:banachmaz} is dense too. Density and category answer different questions. You can approximate any continuous function by polynomials, and a typical continuous function still has no derivative. Volume~12 meets the same split for a random path: a Brownian path is nowhere differentiable with probability one. That statement is a preview.
% MATHFIX: nowhere differentiability of Brownian paths is a preview of Volume~12. It is not proved here and it is not a premise.
\end{beauty}

\section{Points of continuity}

Which subsets of the domain can be exactly the set of continuity points of a function? The oscillation gives the test. For $f:X\to\R$ on a metric space, and for $x\in X$, set
\[\omega_f(x)=\inf_{\delta>0}\ \sup\{|f(y)-f(z)|:\ y,z\in B(x,\delta)\}\in[0,+\infty].\]
The number $\omega_f(x)$ is small when some ball about $x$ has $f$ almost constant on it. Continuity of $f$ at $x$ is exactly $\omega_f(x)=0$. The set $\{\omega_f<\varepsilon\}$ is open. If $\omega_f(x)<\varepsilon$, some ball $B(x,\delta)$ has $|f(y)-f(z)|<\varepsilon$ for all of its points. Any $x'$ in that ball has a smaller ball still inside $B(x,\delta)$, so $\omega_f(x')<\varepsilon$ as well.

\begin{proposition}\label{prop:continuitypoints}
\begin{enumerate}[label=(\alph*)]
\item Let $f:X\to\R$ be any function on a metric space. The set of continuity points of $f$ is a $G_\delta$, namely $\bigcap_k\{\omega_f<\frac1k\}$.
\item The set $\Q$ is not a $G_\delta$ in $\R$. Consequently no function $f:\R\to\R$ is continuous at every rational and discontinuous at every irrational.
\end{enumerate}
\end{proposition}
\begin{proof}
(a) The previous paragraph shows that each set $\{\omega_f<\frac1k\}$ is open, and $f$ is continuous at $x$ exactly when $x$ belongs to all of them.

(b) Suppose $\Q=\bigcap_n U_n$ with each $U_n$ open. Every rational lies in every $U_n$, so each $U_n$ is dense. For each single rational $q$, the complement $\R\setminus\{q\}$ is open and dense. Theorem~\ref{thm:baire} says that the countable family of all the sets $U_n$ and all the sets $\R\setminus\{q\}$ has dense intersection. That intersection is empty: a point of every $U_n$ would be rational, and a point of every $\R\setminus\{q\}$ would be irrational. The empty set is not dense in $\R$.
\end{proof}

The opposite pattern does occur. Thomae's function equals $\frac1q$ at a rational $\frac pq$ in lowest terms, and equals $0$ at every irrational. It is continuous at the irrationals and discontinuous at the rationals. Exercise~\ref{ex:thomae} checks the oscillation at both kinds of point.

\begin{think}
Let $f_n$ be continuous on $[0,1]$, and suppose $f_n(x)\to f(x)$ at every $x$. Can $f$ be discontinuous at every point of $[0,1]$? The indicator of $\Q$ is the natural candidate. Exercise~\ref{ex:v7c05d} asks you to test the formula $\lim_m\lim_n\cos(m!\pi x)^{2n}$ and to see what two successive limits can do. One limit is a different matter, and the next theorem handles it.
\end{think}

\begin{theorem}[Baire, 1899]\label{thm:baireclassone}
Let $X$ be a complete metric space, and let $f_n:X\to\R$ be continuous. Suppose $f_n(x)\to f(x)$ at every $x\in X$. Then the continuity points of $f$ form a comeagre set. In particular $f$ is continuous at the points of a dense set.
\end{theorem}
\begin{proof}
Fix $\varepsilon>0$. For each integer $N\ge1$ define
\[F_N=\{x:|f_m(x)-f_n(x)|\le\varepsilon\text{ for all }m,n\ge N\}.\]
Each condition $|f_m-f_n|\le\varepsilon$ cuts out a closed set, because $f_m$ and $f_n$ are continuous, and $F_N$ is an intersection of those closed sets, so $F_N$ is closed. At each fixed $x$ the numerical sequence $(f_n(x))$ converges, hence is Cauchy, so $x$ belongs to $F_N$ for all large $N$. Thus $\bigcup_N F_N=X$.

Write $U_\varepsilon$ for the union $\bigcup_N\mathring{F_N}$. This set is open, being a union of open sets. It is dense. Take a nonempty open set $G$. Then $G=\bigcup_N(F_N\cap G)$. A nonempty open subset of a complete metric space is not meagre, by Theorem~\ref{thm:baire}, so some piece $F_N\cap G$ fails to be nowhere dense. The closure of that piece is contained in the closed set $F_N$ and has nonempty interior. Pick a point of $F_N\cap G$ that lies in that interior; such a point exists because an open set inside the closure must meet the set itself. Since $G$ is open, a small ball about that point lies in $G$ and in $F_N$. That ball is a nonempty open set $V\subseteq G\cap\mathring{F_N}\subseteq G\cap U_\varepsilon$. Every nonempty open $G$ meets $U_\varepsilon$, so $U_\varepsilon$ is dense.

Now take $x\in\mathring{F_N}$, and choose a ball $B(x,\delta)\subseteq F_N$. For every $y$ in the ball and every $m\ge N$, the inequality $|f_m(y)-f_N(y)|\le\varepsilon$ holds. Let $m\to\infty$. Continuity of the absolute value gives $|f(y)-f_N(y)|\le\varepsilon$. The function $f_N$ is continuous, so a smaller radius makes $|f_N(y)-f_N(z)|<\varepsilon$ for all $y,z$ in the ball. The triangle inequality then yields
\[|f(y)-f(z)|\le|f(y)-f_N(y)|+|f_N(y)-f_N(z)|+|f_N(z)-f(z)|\le3\varepsilon.\]
Hence $\omega_f(x)\le3\varepsilon$ at every point of $U_\varepsilon$. Intersecting over $\varepsilon=\frac1k$ gives
\[\{\omega_f=0\}\supseteq\bigcap_k U_{1/k}.\]
The right-hand side is a countable intersection of dense open sets, so it is a dense $G_\delta$, and it is comeagre by Corollary~\ref{cor:bairecor}(b).
\end{proof}

\begin{corollary}\label{cor:derivative}
If $f:\R\to\R$ is differentiable, then $f'$ is continuous on a comeagre set.
\end{corollary}
\begin{proof}
The difference quotients $x\mapsto n\bigl(f(x+\frac1n)-f(x)\bigr)$ are continuous, and they converge pointwise to $f'(x)$. Theorem~\ref{thm:baireclassone} applies to that sequence.
\end{proof}

The question in the box now has an answer. A single pointwise limit of continuous functions cannot be discontinuous everywhere. In particular $f'$ cannot be the indicator of $\Q$, and it cannot be the indicator of any set whose indicator is discontinuous at every point. A derivative may still fail to be continuous at some points. Extend $x^2\sin\frac1x$ by the value $0$ at $0$. For $x\ne0$ the usual rules give the derivative $2x\sin\frac1x-\cos\frac1x$, and the difference quotient at $0$ is $h\sin\frac1h$, which tends to $0$. The term $-\cos\frac1x$ keeps the derivative from having a limit at $0$.

\section{Category and measure}\label{sec:categorymeasure}

A set of reals can be small in Baire's sense, or small in Lebesgue's sense. The two ideas appeared within a few years of each other, around 1900. They do not describe the same sets. The line is a union of two pieces, one meagre and one of measure zero. John Oxtoby wrote \emph{Measure and Category} about this split. The next chapter starts measure theory properly. One definition is enough for the comparison below.

\begin{definition}
A set $A\subseteq\R$ is \emph{null}, or of \emph{measure zero}, when for every $\varepsilon>0$ the set $A$ can be covered by countably many intervals whose lengths add to less than $\varepsilon$.
\end{definition}

A countable set is null: put an interval of length $\varepsilon2^{-k}$ about the $k$-th point. A countable union of null sets is null: give the $k$-th set a cover of total length $\varepsilon2^{-k}$. The Cantor set is null because it sits in $2^n$ intervals of length $3^{-n}$, and $2^n\cdot3^{-n}$ tends to $0$. An interval $[a,b]$ with $a<b$ is not null. By Theorem~\ref{thm:heineborel}, a cover of $[a,b]$ by open intervals has a finite subcover, so it is enough to show that finitely many open intervals covering $[a,b]$ have total length at least $b-a$. Proceed by induction on the number of intervals. Let $I$ be one interval of the cover that contains $a$, and let $c$ be the right endpoint of $I$. Then $a<c$, the length of $I$ is at least $c-a$, and $I$ does not meet $(c,b]$. The remaining intervals cover $[c,b]$. If $c\ge b$ there is nothing left to cover and the length of $I$ is already at least $b-a$. If $c<b$, the inductive hypothesis gives the remaining intervals total length at least $b-c$, and the sum is at least $b-a$.

\begin{definition}
A real number $x$ is a \emph{Liouville number} when $x$ is irrational and, for every integer $n\ge1$, there exist integers $p$ and $q\ge2$ with
\[\Bigl|x-\frac pq\Bigr|<\frac1{q^n}.\]
\end{definition}

Liouville's constant $\sum_{k\ge1}10^{-k!}$ begins $0.110001000\ldots$, with ones in the decimal places $1!$, $2!$, $3!$, and so on. It is a Liouville number. Take $q=10^{n!}$ and let $p$ be the integer that keeps the digits through place $n!$. The tail is the series from $k=n+1$ onward, and
\[0<\sum_{k\ge n+1}10^{-k!}<2\cdot10^{-(n+1)!}.\]
The comparison $(n+1)!\ge n\cdot n!+n!$ gives $2\cdot10^{-(n+1)!}\le q^{-n}$. For a numerical check, take $n=1$. Then $q=10$, and the tail after the first decimal digit is $0.010001000\ldots$, which is less than $\frac1{10}$. Liouville numbers are transcendental. Exercise~\ref{ex:liouville} proves that, and so does Volume~2, Exercise~9.7. They were the first transcendental numbers written down by an explicit rule.

\begin{theorem}\label{thm:liouville}
Let $\mathcal L$ be the set of Liouville numbers. Then $\mathcal L$ is a dense $G_\delta$ in $\R$, hence comeagre, and $\mathcal L$ is null. The line therefore splits as the disjoint union $\R=\mathcal L\cup(\R\setminus\mathcal L)$ of a null set and a meagre set.
\end{theorem}
\begin{proof}
For each $n\ge1$ define
\[U_n=\bigcup_{q\ge2}\bigcup_{p\in\Z}\,\Bigl]\frac pq-q^{-n},\frac pq+q^{-n}\Bigr[.\]
Each $U_n$ is a union of open intervals, so $U_n$ is open. Every rational $\frac pq$ with $q\ge2$ lies in $U_n$, so $U_n$ is dense. A Liouville number lies in every $U_n$ and in none of the singletons $\{r\}$ with $r\in\Q$. Conversely, an irrational point of every $U_n$ satisfies the Liouville inequalities. Thus
\[\mathcal L=\bigcap_n U_n\cap\bigcap_{r\in\Q}(\R\setminus\{r\}).\]
This is a countable intersection of dense open sets, so it is a dense $G_\delta$. Corollary~\ref{cor:bairecor}(b) says that a dense $G_\delta$ is comeagre.

For the measure, it is enough to show that each piece $\mathcal L\cap[-m,m]$ is null, because $\mathcal L$ is the union of those pieces over integers $m\ge1$. Fix $m\ge1$ and $n\ge3$. The piece sits in the intervals $\bigl]\frac pq-q^{-n},\frac pq+q^{-n}\bigr[$ with $q\ge2$ and $|p|\le mq+1$. For each $q$ the number of such integers $p$ is at most $2mq+3$. For $m\ge1$ and $q\ge2$ one has $2mq+3\le5mq$. Each interval has length $2q^{-n}$, so the total length is at most
\[\sum_{q\ge2}5mq\cdot\frac2{q^n}=10m\sum_{q\ge2}\frac1{q^{n-1}}\le10m\int_1^\infty\frac{\dd t}{t^{n-1}}=\frac{10m}{n-2}.\]
The integral comparison is the usual one for a decreasing function: the box of height $q^{-(n-1)}$ on $[q-1,q]$ sits under the graph of $t\mapsto t^{-(n-1)}$ on $[1,\infty)$. The bound $\frac{10m}{n-2}$ tends to $0$ as $n\to\infty$. Hence $\mathcal L\cap[-m,m]$ is null, and so is $\mathcal L$. The complement of a comeagre set is meagre, so $\R\setminus\mathcal L$ is meagre.
\end{proof}

\begin{think}
Which reading of ``most'' do you trust for rational approximation? Is a typical real number easy to approximate by fractions, or hard? And is a typical number normal?
\end{think}

The answer depends on the reading. In the sense of category, a typical real is Liouville, by Theorem~\ref{thm:liouville}: for every $n$ some fraction $\frac pq$ lands inside $q^{-n}$ of it. In the sense of measure, almost no real is Liouville, because $\mathcal L$ is null. Khinchin's theorem of 1924 gives the scale that measure actually sees. For almost every real $x$, the inequality $\bigl|x-\frac pq\bigr|<\frac1{q^2\log q}$ has infinitely many solutions in integers $p$ and $q\ge2$, while $\bigl|x-\frac pq\bigr|<\frac1{q^2(\log q)^2}$ has only finitely many. The proof is in Chapter~\ref{ch:harder}. It is not used as a premise in this chapter.
% MATHFIX: Khinchin's theorem is a preview of Chapter~\ref{ch:harder}. The inequalities are quoted here and are not proved in this chapter.
Normality splits the same way. Almost every number is normal. That is Borel's theorem, proved in Chapter~\ref{ch:measures}, and it is not a premise here.
% MATHFIX: Borel's normality theorem is a preview of Chapter~\ref{ch:measures}. It is not proved in this chapter.
A typical number, in Baire's sense, fails to be normal. The next proposition proves that half.

\begin{proposition}\label{prop:nonnormal}
For $x\in[0,1]$, let $\nu_n(x)$ be the proportion of ones among the first $n$ binary digits of $x$, using the greedy expansion of Proposition~\ref{prop:cantor}. The set of $x\in[0,1]$ with $\limsup_n\nu_n(x)=1$ is comeagre in $[0,1]$. The binary normal numbers in $[0,1]$ therefore form a meagre set.
\end{proposition}
\begin{proof}
A word $w$ of length $n$ names the half-open cylinder of numbers whose greedy expansion begins with $w$. The corresponding open interval is $I_w=\,]0.w,\,0.w+2^{-n}[$, of length $2^{-n}$. For each integer $N\ge1$, let $V_N$ be the union of those $I_w$ for which the length $n$ of $w$ is at least $N$ and $w$ contains more than $\bigl(1-\frac1N\bigr)n$ ones. Each $V_N$ is open. It is dense in $[0,1]$. Any nonempty open interval in $[0,1]$ contains some cylinder $I_u$. Append enough ones to $u$ to produce a word $w$ of length at least $N$ whose proportion of ones exceeds $1-\frac1N$. Then $I_w\subseteq I_u$, so the original interval meets $V_N$. The interval $[0,1]$ is complete. Theorem~\ref{thm:baire} says that $\bigcap_N V_N$ is a dense $G_\delta$, hence comeagre by Corollary~\ref{cor:bairecor}(b). Take $x$ in that intersection. For every $N$ there is some $n\ge N$ with $\nu_n(x)>1-\frac1N$, so $\limsup_n\nu_n(x)=1$. A binary normal number has $\nu_n(x)\to\frac12$, so it lies outside the intersection.
\end{proof}

\begin{remark}
The parallel runs deeper than these two examples, once one assumes the continuum hypothesis. Sierpi\'nski and Erd\H{o}s proved, in work published from 1934 to 1943, that there is a bijection $f:\R\to\R$ with $f^{-1}=f$ which sends meagre sets to null sets and null sets to meagre sets. A theorem about meagre sets that uses only the abstract properties shared with null sets then has a twin about null sets. Oxtoby's \emph{Measure and Category} is the account of that correspondence.
% MATHFIX: the Sierpiński--Erdős bijection, under the continuum hypothesis, is quoted in this remark and is not proved in this chapter.
\end{remark}

\section{What we proved, and what comes next}

Theorem~\ref{thm:baire} is the sentence the chapter keeps using. In a complete metric space, or in a locally compact Hausdorff space, a countable intersection of dense open sets is still dense. A property is typical when it holds on a dense $G_\delta$, and countably many typical properties hold at the same typical points.

The three questions in the opening now have category answers. A typical function in $C([0,1])$ has a derivative at no point, by Theorem~\ref{thm:banachmaz}. Takagi's series is one function you can write down and test by hand, by Theorem~\ref{thm:takagi}. A typical real number is Liouville, so it is approximated by rationals more closely than any fixed power of the denominator suggests, by Theorem~\ref{thm:liouville}. A typical binary expansion is not normal, by Proposition~\ref{prop:nonnormal}. Along the way, Theorem~\ref{thm:baireclassone} showed that a pointwise limit of continuous functions is continuous on a comeagre set, so a derivative cannot be discontinuous at every point.

Measure reads two of those three questions the other way. The Liouville numbers are null. Chapter~\ref{ch:measures} will prove that almost every number is normal. The line is a union of a meagre set and a null set, by Theorem~\ref{thm:liouville}. Each notion of small is doing its own job. From Chapter~\ref{ch:measures} on, the volume builds the measure-theoretic one.

The next chapter starts from a fact that is easy to believe and harder to prove. A closed curve in the plane that never crosses itself should separate an inside from an outside. The curves in this chapter, some of them without a tangent at any point, make that belief something one ought to prove. Chapter~\ref{ch:jordan} proves it.

\section{Exercises}

\begin{exercise}[$\star$]\label{ex:thomae}
\emph{Thomae's function.} Define $f(x)=\frac1q$ when $x=\frac pq$ in lowest terms with $q\ge1$, and define $f(x)=0$ when $x$ is irrational. Prove that $f$ is continuous at every irrational and discontinuous at every rational.
\end{exercise}

\begin{exercise}[$\star$]\label{ex:perfect}
Let $X$ be a nonempty complete metric space with no isolated point. Prove that $X$ is uncountable. Deduce, again, that $\R$, the Cantor set, and $\Z_p$ (Chapter~\ref{ch:metric}) are uncountable.
\end{exercise}

\begin{exercise}[$\star\star$]\label{ex:hamelbanach}
Let $E$ be a Banach space whose dimension is infinite. Prove that any vector-space basis of $E$, in the sense of Theorem~\ref{thm:basis}, is uncountable. A finite-dimensional subspace is complete, hence closed, by Corollary~\ref{cor:norms}. Deduce that no norm on the space of polynomials is complete.
\end{exercise}

\begin{exercise}[$\star\star$]\label{ex:monotone2}
Prove that a typical $f\in C([0,1])$ fails to be monotone on every non-degenerate subinterval of $[0,1]$. For an interval $I$ with rational endpoints, consider the functions that are non-decreasing on $I$.
\end{exercise}

\begin{exercise}[$\star\star$]\label{ex:liouville}
\emph{Liouville's theorem.} Let $x$ be a real root of an irreducible polynomial $P\in\Z[X]$ of degree $d\ge2$. Prove that some constant $c>0$ satisfies $\bigl|x-\frac pq\bigr|\ge\frac c{q^d}$ for all integers $p$ and all integers $q\ge1$. Deduce that every Liouville number is transcendental.
\end{exercise}

\begin{exercise}[$\star\star$]\label{ex:fsigma}
Let $E\subseteq\R$ be an $F_\sigma$ set. Construct a function $f:\R\to\R$ whose set of discontinuities is exactly $E$. Write $E=\bigcup_n F_n$ with the closed sets $F_n$ increasing, and combine scaled indicators of $F_n\cap\Q$ and of $F_n\setminus\mathring{F_n}$. Together with Proposition~\ref{prop:continuitypoints}(a), the construction characterises the sets that arise as discontinuity sets.
\end{exercise}

\begin{exercise}[$\star\star\star$]\label{ex:polynomial}
Let $f:\R\to\R$ be infinitely differentiable, and suppose that at every $x$ some derivative vanishes: $f^{(n)}(x)=0$ for an integer $n=n(x)$. Prove that $f$ is a polynomial. If a single $n$ worked at every $x$, the claim would be immediate. Baire's theorem produces an open piece on which one $n$ does work.
\end{exercise}

\begin{exercise}[$\star$]\label{ex:takagivalues}
\emph{A computation with Takagi's function.} (a) Compute $T\bigl(\frac12\bigr)$, $T\bigl(\frac14\bigr)$, and $T\bigl(\frac13\bigr)$. (b) Prove that $\phi(y)+\frac12\phi(2y)\le\frac12$ for every real $y$, and deduce that $T(x)\le\frac23$ for every $x$. The value $\frac13$ is a point where the bound is attained. (c) For $x=\frac13$, compute the chord slopes $s_1,s_2,s_3,s_4$ from the proof of Theorem~\ref{thm:takagi}.
\end{exercise}

\begin{exercise}[$\star$]\label{ex:v7c05a}
Prove that a finite union of nowhere dense sets is nowhere dense. Conclude that every finite subset of $\R$ is nowhere dense, and that $\Z$ is nowhere dense in $\R$. Then show that a countable union of nowhere dense sets can be dense: $\Q$ is meagre in $\R$, and $\Q$ is not nowhere dense.
\end{exercise}

\begin{exercise}[$\star\star$]\label{ex:v7c05b}
Let $P_n$ be the set of polynomial functions of degree at most $n$, viewed as a subset of $C([0,1])$ with the supremum norm. Prove that each $P_n$ is closed and has empty interior. Deduce that the polynomials form a meagre subset of $C([0,1])$. They are also dense, by Corollary~\ref{cor:weierstrass}. A typical continuous function is therefore not a polynomial, while polynomials still come arbitrarily close to every continuous function.
\end{exercise}

\begin{exercise}[$\star\star$]\label{ex:v7c05c}
Prove that a typical $f\in C([0,1])$ does not vanish identically on any non-degenerate closed subinterval of $[0,1]$. For rational endpoints $a<b$ in $[0,1]$, consider the set of functions that are identically zero on $[a,b]$.
\end{exercise}

\begin{exercise}[$\star\star\star$]\label{ex:v7c05d}
For integers $m\ge1$ and $n\ge1$, set $f_{m,n}(x)=\cos(m!\pi x)^{2n}$. (a) Show that $g_m(x)=\lim_n f_{m,n}(x)$ exists for every real $x$, and that $g_m(x)=1$ when $x$ is an integer multiple of $\frac1{m!}$, while $g_m(x)=0$ otherwise. (b) Show that $\lim_m g_m(x)$ equals $1$ when $x$ is rational and equals $0$ when $x$ is irrational. (c) Deduce that the indicator of $\Q$ is discontinuous at every real, and that this indicator is not a pointwise limit of continuous functions on $\R$.
\end{exercise}

\section{Hints}

\begin{hint}{ex:thomae}
Around an irrational, only finitely many fractions with denominator at most $N$ lie in a bounded interval, and the irrational equals none of them. Around a rational, the value of $f$ is positive and the irrational neighbours have $f=0$.
\end{hint}
\begin{hint}{ex:perfect}
A singleton is closed. With no isolated points it is also nowhere dense. Theorem~\ref{thm:baire} forbids a nonempty complete space to be a countable union of such singletons.
\end{hint}
\begin{hint}{ex:hamelbanach}
A countable basis would write $E$ as the union of the spans $E_n=\operatorname{span}(e_1,\dots,e_n)$. A proper subspace of a normed space has empty interior: a ball inside it can be translated and scaled to fill the whole space.
\end{hint}
\begin{hint}{ex:monotone2}
The set of functions that are non-decreasing on a fixed $[a,b]$ is closed under uniform limits. Near any $f$, a polynomial plus a steep sawtooth decreases on a small piece inside $[a,b]$. Countably many rational endpoints see every subinterval.
\end{hint}
\begin{hint}{ex:liouville}
Irreducibility keeps $P\bigl(\frac pq\bigr)$ off zero, and $q^d P\bigl(\frac pq\bigr)$ is a nonzero integer, so $\bigl|P\bigl(\frac pq\bigr)\bigr|\ge q^{-d}$. Compare that lower bound with the mean value theorem on $[x-1,x+1]$. Volume~1, Proposition~6.9, is Gauss's lemma, if you want the factorisation form of the rational-root step.
\end{hint}
\begin{hint}{ex:fsigma}
Inside the interior of $F_n$, the indicator of the rationals jumps at every point. On the boundary $F_n\setminus\mathring{F_n}$, the indicator of $F_n$ itself jumps. Weights $3^{-n}$ keep the first nonzero jump visible above the tail. Outside $E$, only a tail remains, and the tail is small.
\end{hint}
\begin{hint}{ex:polynomial}
Let $G$ be the union of the open intervals on which $f$ agrees with a polynomial, and set $F=\R\setminus G$. Show that $F$ has no isolated point. Then apply Corollary~\ref{cor:bairecor}(a) inside $F$ to the closed sets $\{x\in F:f^{(n)}(x)=0\}$.
\end{hint}

\begin{hint}{ex:takagivalues}
(a) At a dyadic rational the series stops after finitely many terms. At $\frac13$, the number $2^n/3$ is always at distance $\frac13$ from the nearest integer. (b) Period $1$ and the symmetry $y\mapsto1-y$ reduce the inequality to $y\in[0,\frac12]$. Group the series for $T$ in pairs. (c) The sign $\varepsilon_n$ equals $+1$ when the fractional part of $2^nx$ is strictly below $\frac12$.
\end{hint}

\begin{hint}{ex:v7c05a}
The closure of a union of two sets is the union of the closures. If an open set sat in that union, deleting one closed piece would leave a nonempty open set in the other, or else the first piece would already have interior.
\end{hint}
\begin{hint}{ex:v7c05b}
Corollary~\ref{cor:norms} makes a finite-dimensional normed space complete, and Proposition~\ref{prop:complete}(a) turns completeness into closedness. If a ball sat inside $P_n$, translation and scaling would force $P_n=C([0,1])$. The function $|x-\frac12|$ is continuous and is not a polynomial.
\end{hint}
\begin{hint}{ex:v7c05c}
For rational $a<b$, the set of functions vanishing on $[a,b]$ is closed. At least one of $f+\frac\varepsilon2$ and $f-\frac\varepsilon2$ fails to vanish on the whole of $[a,b]$, because their difference is the constant $\varepsilon$. Every non-degenerate subinterval contains one with rational endpoints.
\end{hint}
\begin{hint}{ex:v7c05d}
The power $2n$ sends a number of absolute value strictly less than $1$ to $0$, and it leaves $\pm1$ at the value $1$. The equality $|\cos(m!\pi x)|=1$ means that $m!\,x$ is an integer. For a rational $a/b$, that happens as soon as $m\ge b$. Theorem~\ref{thm:baireclassone} then forbids the indicator of $\Q$ to be a single pointwise limit of continuous functions.
\end{hint}

\section{Solutions}

\begin{solution}{ex:thomae}
\idea At an irrational, fractions with a small denominator are isolated, so a small enough neighbourhood keeps them out and forces $f$ below any given $\varepsilon$. At a rational the function is positive, while every neighbourhood contains irrationals, where the function is $0$.
\writeup Let $x=\frac pq$ in lowest terms, so $f(x)=\frac1q>0$. Any sequence of irrationals tending to $x$ has $f=0$ along the sequence. The function values do not tend to $f(x)$, so $f$ is discontinuous at $x$.

Now let $x$ be irrational, and fix $\varepsilon>0$. Choose an integer $N$ with $\frac1N<\varepsilon$. In the interval $[x-1,x+1]$ there are only finitely many fractions $\frac pq$ in lowest terms with $q\le N$: for each fixed $q$ there are at most $2q+1$ candidates. None of them equals $x$. The distance from $x$ to the nearest one is positive, so some $\delta\in\,]0,1[$ makes $]x-\delta,x+\delta[$ avoid all of them. On that neighbourhood, an irrational $y$ has $f(y)=0$, and a rational $y=\frac pq$ in lowest terms has $q>N$, hence $f(y)=\frac1q<\frac1N$. Since $f(x)=0$, one gets $|f(y)-f(x)|=f(y)<\varepsilon$. That estimate is continuity of $f$ at this irrational $x$.
\end{solution}

\begin{solution}{ex:perfect}
\idea Countability would write $X$ as a union of singletons. A space with no isolated point makes each singleton nowhere dense, and Theorem~\ref{thm:baire} says a nonempty complete space is not a countable union of nowhere dense sets.
\writeup The singleton $\{x\}$ is closed. It is not open: an open singleton would be an isolated point. So $\{x\}$ is nowhere dense. If $X$ were countable it would be meagre in itself, contradicting Theorem~\ref{thm:baire}.

The line $\R$ is complete and has no isolated point. The Cantor set is closed in $\R$, hence complete by Proposition~\ref{prop:complete}(a), and it has no isolated point. If $x=h(a)$ in the coding of Proposition~\ref{prop:cantor}, and if $a'$ differs from $a$ only in the $n$-th place, then $|h(a)-h(a')|=2\cdot3^{-n}$. Those neighbours are distinct from $x$ and approach $x$ as $n$ grows. The group $\Z_p$ is complete, and $x+p^n$ lies at distance $p^{-n}$ from $x$, so no point of $\Z_p$ is isolated.
\end{solution}

\begin{solution}{ex:hamelbanach}
\idea A countable basis would present the Banach space as a countable union of finite-dimensional subspaces. Each of those is closed and has empty interior. Theorem~\ref{thm:baire} does not allow that union to be the whole space.
\writeup A finite basis is excluded by the hypothesis that $E$ is infinite-dimensional. Suppose $(e_k)_{k\ge1}$ is a countable basis, and set $E_n=\operatorname{span}(e_1,\dots,e_n)$, so that $E=\bigcup_n E_n$. The subspace $E_n$ is finite-dimensional. Corollary~\ref{cor:norms} says that the restricted norm makes $E_n$ complete, and Proposition~\ref{prop:complete}(a) says that a complete subset is closed, so $E_n$ is closed in $E$.

The interior is empty. If some ball $B(x,r)$ sat inside $E_n$, then $B(0,r)=B(x,r)-x$ would sit inside $E_n$ as well, because $E_n$ is a subspace. For $y\ne0$ the vector $\frac{r}{2\|y\|}y$ has norm $\frac r2$, so it lies in $B(0,r)$ and therefore in $E_n$. Scaling back puts $y$ in $E_n$. Thus $E_n=E$. That is impossible, because $e_{n+1}\notin E_n$. Each $E_n$ is nowhere dense, the union is all of $E$, and Theorem~\ref{thm:baire} is contradicted.

The monomials $(x^k)_{k\ge0}$ are a countable basis of the polynomials. The previous paragraph says that no norm can make that space complete.
\end{solution}

\begin{solution}{ex:monotone2}
\idea Monotonicity on one fixed interval is a closed condition, and a steep sawtooth destroys it inside every ball. Rational endpoints reduce the uncountable family of all intervals to a countable family, which is what Baire can absorb.
\writeup Let $a<b$ be rational numbers in $[0,1]$. Write $M^+_{a,b}$ for the set of $f\in C([0,1])$ that are non-decreasing on $[a,b]$, and $M^-_{a,b}$ for the non-increasing ones. If $f_k\in M^+_{a,b}$ and $f_k\to f$ uniformly, the inequalities $f_k(x)\le f_k(y)$ for $a\le x\le y\le b$ pass to the limit, so $M^+_{a,b}$ is closed. The same holds for $M^-_{a,b}$.

The interior is empty. Fix $f$ and $\varepsilon>0$. Choose a polynomial $p$ with $\|f-p\|_\infty<\frac\varepsilon2$, and let $M$ be large enough that $p$ is $M$-Lipschitz on $[0,1]$, using Theorem~\ref{thm:mvt} as in the proof of Theorem~\ref{thm:banachmaz}. Inside $[a,b]$, build a sawtooth of height at most $\frac\varepsilon2$ that has one piece of slope $-L$ with $L>M$. The sum $g=p+s$ satisfies $\|g-f\|_\infty<\varepsilon$, and on that piece the slope of $g$ is at most $-L+M<0$, so $g$ is not non-decreasing on $[a,b]$. The same construction with slope $+L$ handles $M^-_{a,b}$.

Suppose $f$ is monotone on some non-degenerate subinterval of $[0,1]$. That subinterval contains a smaller interval $[a,b]$ with rational endpoints, and $f$ is monotone on $[a,b]$ as well. Thus $f$ lies in $\bigcup(M^+_{a,b}\cup M^-_{a,b})$, the union running over rational pairs $a<b$ in $[0,1]$. Each term is nowhere dense, the index set is countable, and the union is meagre. Theorem~\ref{thm:baire} says that a typical $f$ lies outside the union.
\end{solution}

\begin{solution}{ex:liouville}
\idea An algebraic number of degree $d$ stays at least a constant times $q^{-d}$ away from every fraction $\frac pq$. The reason is that $q^d P\bigl(\frac pq\bigr)$ is a nonzero integer, while $P(x)=0$ and $P$ is Lipschitz near $x$. A Liouville number violates every such lower bound.
\writeup Let $M=\max_{[x-1,x+1]}|P'|$, which is finite because $P'$ is continuous. The degree is at least $2$, so $P'$ is not the zero polynomial, and $M>0$. Set $c=\min\bigl(1,\frac1M\bigr)$.

Take integers $p$ and $q\ge1$. If $\bigl|x-\frac pq\bigr|\ge1$, then $\bigl|x-\frac pq\bigr|\ge c$, and $q^d\ge1$ gives $\bigl|x-\frac pq\bigr|\ge\frac c{q^d}$. Now suppose $\bigl|x-\frac pq\bigr|<1$. The value $P\bigl(\frac pq\bigr)$ is not zero. Indeed $P$ is irreducible of degree $d\ge2$, so it has no rational root: a rational root would give a linear factor over $\Q$, and Volume~1, Proposition~6.9, would turn that factorisation into a non-unit factorisation in $\Z[X]$. Clearing the denominator, $q^d P\bigl(\frac pq\bigr)$ is an integer, and it is nonzero, so its absolute value is at least $1$. Therefore $\bigl|P\bigl(\frac pq\bigr)\bigr|\ge q^{-d}$. The mean value theorem, Theorem~\ref{thm:mvt}, supplies a point between $x$ and $\frac pq$, hence in $[x-1,x+1]$, with
\[\frac1{q^d}\le\Bigl|P\Bigl(\frac pq\Bigr)-P(x)\Bigr|\le M\Bigl|\frac pq-x\Bigr|.\]
Thus $\bigl|x-\frac pq\bigr|\ge\frac1{Mq^d}\ge\frac c{q^d}$.

Let $x$ be Liouville, and suppose toward a contradiction that $x$ were algebraic. Its minimal polynomial may be taken in $\Z[X]$, irreducible, of some degree $d\ge2$, because a Liouville number is irrational. The previous paragraph gives a constant $c>0$. For any $n>d$ the Liouville condition produces $q\ge2$ and an integer $p$ with $\frac c{q^d}\le\bigl|x-\frac pq\bigr|<\frac1{q^n}$. Then $c<q^{d-n}\le2^{d-n}$. The right-hand side tends to $0$ as $n$ grows, so the inequality fails for large $n$.
\end{solution}

\begin{solution}{ex:fsigma}
\idea A closed set $F$ produces discontinuities of two kinds. In the interior, the indicator of the rationals jumps everywhere. On the boundary, the indicator of $F$ jumps. Fast-decaying weights $3^{-n}$ let the earliest jump dominate the tail, and outside the union every weight eventually vanishes in a neighbourhood.
\writeup Replace $F_n$ by $F_1\cup\cdots\cup F_n$, so that $E=\bigcup_n F_n$ with $F_1\subseteq F_2\subseteq\cdots$ and each $F_n$ closed. Define
\[g_n=\mathbf 1_{F_n}+\mathbf 1_{F_n\cap\Q}\in\{0,1,2\},\qquad f=\sum_{n\ge1}3^{-n}g_n.\]
The series converges uniformly, because $0\le g_n\le2$.

\emph{Continuity off $E$.} If $x\notin E$, then $x\notin F_n$ for every $n$. Each $F_n$ is closed, so $g_n$ vanishes on a neighbourhood of $x$. Fix $\varepsilon>0$ and choose $N$ with $\sum_{n>N}2\cdot3^{-n}<\varepsilon$. Shrink the neighbourhood so that $g_1,\dots,g_N$ all vanish on it. There $f(x)=0$ and $0\le f<\varepsilon$, so $f$ is continuous at $x$.

\emph{Discontinuity on $E$.} Let $x\in E$, and let $n$ be the least index with $x\in F_n$. The closed sets $F_m$ for $m<n$ miss $x$, so some neighbourhood $V$ of $x$ kills $g_m$ for all $m<n$.

Suppose first that $x\in\mathring{F_n}$. Shrink $V$ into $F_n$. The sets are nested, so $V\subseteq F_m$ for every $m\ge n$. On $V$,
\[f(y)=c(y)\sum_{m\ge n}3^{-m}=\frac32\,c(y)\,3^{-n},\]
where $c(y)=2$ if $y$ is rational and $c(y)=1$ if $y$ is irrational. The two values are $3^{1-n}$ and $\frac32\cdot3^{-n}$. Every neighbourhood of $x$ contains both rationals and irrationals, so both values occur arbitrarily close to $x$, and $f$ is discontinuous at $x$.

Suppose instead that $x\in F_n\setminus\mathring{F_n}$. Points $y\in V\setminus F_n$ come arbitrarily close to $x$. For such a $y$ one has $g_m(y)=0$ for all $m\le n$, hence
\[f(y)\le\sum_{m>n}2\cdot3^{-m}=3^{-n}.\]
At $x$ itself, $g_m(x)\ge1$ for every $m\ge n$, so
\[f(x)\ge\sum_{m\ge n}3^{-m}=\frac32\cdot3^{-n}.\]
Therefore $|f(y)-f(x)|\ge\frac12\cdot3^{-n}$ for those nearby $y$, and $f$ is discontinuous at $x$.

Proposition~\ref{prop:continuitypoints}(a) says that every discontinuity set is a $G_\delta$ complement, hence an $F_\sigma$. The construction shows that every $F_\sigma$ arises. The two statements together are the characterisation.
\end{solution}

\begin{solution}{ex:polynomial}
\idea Wherever the vanishing order $n(x)$ is bounded on a whole interval, $f$ agrees with a polynomial there. Baire's theorem, applied inside the closed set where $f$ is not locally a polynomial, produces an interval on which one derivative vanishes throughout that set. The gaps of the set inherit the same vanishing, and the interval was not supposed to be good.
\writeup Let $G$ be the union of all open intervals on which $f$ coincides with some polynomial, and set $F=\R\setminus G$. The set $G$ is open, so $F$ is closed.

\emph{Two facts.} (i) Suppose $f^{(m)}(x_k)=0$ for a sequence $x_k\to x$ with $x_k\ne x$. Then $f^{(j)}(x)=0$ for every $j\ge m$. Passing to a subsequence, one may assume that $(x_k)$ is strictly monotone. Take the decreasing case; the increasing case is the same with the intervals reversed. Continuity of $f^{(m)}$ gives $f^{(m)}(x)=0$. On each interval between $x_{k+1}$ and $x_k$, Theorem~\ref{thm:mvt} applied to $f^{(m)}$ produces a point $x'_k$ strictly between them with $f^{(m+1)}(x'_k)=0$. The points $x'_k$ are distinct from $x$ and tend to $x$. Run the same argument with $m+1$ in place of $m$, and the higher derivatives vanish at $x$ as well. (ii) If $f$ equals a polynomial $p$ on an open interval $]a,c[$, continuity of each derivative gives $f^{(j)}(c)=p^{(j)}(c)$ for every $j$. The same holds at a left endpoint. So the derivatives of $f$ at an endpoint determine the polynomial.

\emph{$F$ has no isolated point.} If $x\in F$ were isolated, then $f$ would coincide with a polynomial $p$ on some $]a,x[$ and with a polynomial $q$ on some $]x,b[$. Fact~(ii) says that every derivative of $p$ at $x$, and every derivative of $q$ at $x$, equals the corresponding derivative of $f$. Two polynomials with the same derivatives of all orders at one point are equal: their difference has Taylor series zero at that point. Thus $p=q$. Continuity of $f$ at $x$ gives $f=p$ on $]a,b[$, so $x\in G$, which contradicts $x\in F$.

\emph{The set $F$ has to be empty.} Suppose, on the contrary, that it is not. The set $F$ is closed in $\R$, hence complete. By hypothesis $F=\bigcup_n\{x\in F:f^{(n)}(x)=0\}$, and each of those sets is closed in $F$ because $f^{(n)}$ is continuous. Corollary~\ref{cor:bairecor}(a), applied in $F$, says that some piece has nonempty interior relative to $F$. So there exist an integer $n$ and an open interval $J=\,]a,b[$ with $J\cap F\ne\varnothing$ and $f^{(n)}=0$ on $J\cap F$. No point of $F$ is isolated, so each $x\in J\cap F$ is a limit of other points of $J\cap F$. Fact~(i) upgrades the vanishing: $f^{(j)}(x)=0$ for every $j\ge n$ and every $x\in J\cap F$.

Let $]c,d[$ be a connected component of $J\cap G$, so $]c,d[$ is a maximal open interval contained in $J\cap G$. The interval $J$ is not contained in $G$, because $J$ meets $F$. One endpoint of $]c,d[$ therefore lies in $J\cap F$. The two ends are symmetric; take the left end $c\in J\cap F$. On $]c,d[$ the function $f$ is locally equal to a polynomial. Fact~(ii) says that these local polynomials agree on overlaps, so a single polynomial $p$ works on all of $]c,d[$. In more detail, fix one local representation $f=p$ on a subinterval, and let $S$ be the set of points of $]c,d[$ near which $f=p$. Then $S$ is open in $]c,d[$, and fact~(ii) makes $S$ closed in $]c,d[$, so $S=\,]c,d[$. Applying fact~(ii) at the endpoint $c$ gives $p^{(j)}(c)=f^{(j)}(c)=0$ for every $j\ge n$. A polynomial whose derivatives of order $n$ and higher vanish at a point has degree less than $n$. Thus $f^{(n)}=0$ on $]c,d[$, that is, on this piece of $J\cap G$. Combined with the vanishing already proved on $J\cap F$, one gets $f^{(n)}=0$ on all of $J$. So $f$ coincides on $J$ with a polynomial of degree less than $n$, which means $J\subseteq G$. That contradicts $J\cap F\ne\varnothing$.

\emph{One polynomial on the line.} The case $F\ne\varnothing$ is impossible, so $G=\R$. Thus $f$ is locally a polynomial on $\R$. The same open-and-closed argument as above, using connectedness of $\R$ and fact~(ii), promotes the local polynomials to one polynomial on the whole line.
\end{solution}

\begin{solution}{ex:takagivalues}
\idea At a dyadic point only finitely many tents are awake. At $\frac13$ every tent has the same height, and the series is geometric. Pairing terms turns the inequality $\phi(y)+\frac12\phi(2y)\le\frac12$ into the bound $T\le\frac23$.
\writeup (a) Since $\phi$ vanishes on $\Z$, one has $T\bigl(\frac12\bigr)=\phi\bigl(\frac12\bigr)=\frac12$, and every later term is $\phi$ at an integer. Next, $T\bigl(\frac14\bigr)=\phi\bigl(\frac14\bigr)+\frac12\phi\bigl(\frac12\bigr)=\frac14+\frac14=\frac12$, and $2^n/4$ is an integer for all $n\ge2$. For the third value, $2^n$ modulo $3$ is $1$ or $2$, so $\frac{2^n}3$ is an integer plus $\frac13$ or plus $\frac23$. In both cases $\phi\bigl(\frac{2^n}3\bigr)=\frac13$. Therefore
\[T\Bigl(\frac13\Bigr)=\frac13\sum_{n\ge0}2^{-n}=\frac23.\]

(b) Both sides of the proposed inequality have period $1$, and both are unchanged when $y$ is replaced by $1-y$. It is enough to treat $y\in[0,\frac12]$, where $\phi(y)=y$. If $0\le y\le\frac14$, then $\phi(2y)=2y$ and the left side equals $y+y=2y\le\frac12$. If $\frac14\le y\le\frac12$, then $\phi(2y)=1-2y$ and the left side equals $y+\frac12(1-2y)=\frac12$. Group the series for $T$ in pairs:
\[T(x)=\sum_{m\ge0}4^{-m}\Bigl(\phi(4^mx)+\frac12\phi(2\cdot4^mx)\Bigr)\le\sum_{m\ge0}4^{-m}\cdot\frac12=\frac23.\]
Part~(a) shows that $T\bigl(\frac13\bigr)=\frac23$, so the constant is attained.

(c) In the proof of Theorem~\ref{thm:takagi}, the sign $\varepsilon_n$ is the slope of $x\mapsto 2^{-n}\phi(2^nx)$ on the dyadic interval that contains $x$. That slope is $+1$ when the fractional part of $2^nx$ lies in $[0,\frac12)$, and it is $-1$ when the fractional part lies in $(\frac12,1)$. For $x=\frac13$ the fractional parts of $x$, $2x$, $4x$, and $8x$ are $\frac13$, $\frac23$, $\frac13$, and $\frac23$. The signs $\varepsilon_0,\varepsilon_1,\varepsilon_2,\varepsilon_3$ are therefore $+1,-1,+1,-1$, and the partial sums are $s_1,s_2,s_3,s_4=1,0,1,0$. The chord slopes over the shrinking dyadic intervals about $\frac13$ oscillate between $1$ and $0$. At the highest point of $T$, the curve still has no tangent.
\end{solution}

\begin{solution}{ex:v7c05a}
\idea Two nowhere dense sets stay nowhere dense because the closure of a union is the union of the closures, and a nonempty open set cannot hide in a union of two closed empty-interior sets. Countable unions are a different matter: the singletons of the rationals are nowhere dense, and their union is dense.
\writeup Let $A$ and $B$ be nowhere dense, and suppose a nonempty open set $U$ sat inside $\overline{A\cup B}=\bar A\cup\bar B$. The difference $U\setminus\bar A$ is open. If it is empty, then $U\subseteq\bar A$, so $\bar A$ has nonempty interior, contradicting the choice of $A$. If it is nonempty, then it is a nonempty open set contained in $\bar B$, and $\bar B$ has nonempty interior. Either way one of the two sets fails to be nowhere dense. So $A\cup B$ is nowhere dense. A finite union follows by induction on the number of sets.

A singleton in $\R$ is closed and contains no open interval, so it is nowhere dense. The previous paragraph puts every finite set in the same class. The integers are closed in $\R$, being the complement of $\bigcup_n\,]n,n+1[$, and $\Z$ contains no open interval, so $\Z$ itself is nowhere dense. In particular a countable union of nowhere dense sets can be nowhere dense.

It can also be dense. Each singleton $\{q\}$ with $q\in\Q$ is nowhere dense, and $\Q$ is their union, so $\Q$ is meagre. The closure of $\Q$ is $\R$, which has nonempty interior, so $\Q$ is not nowhere dense.
\end{solution}

\begin{solution}{ex:v7c05b}
\idea Each $P_n$ is a finite-dimensional subspace of $C([0,1])$. Finite-dimensional normed spaces are complete, complete subsets are closed, and a proper subspace cannot contain a ball. Polynomials of all degrees are still dense, by the Weierstrass theorem already proved.
\writeup The set $P_n$ is a vector subspace of $C([0,1])$ of dimension $n+1$, with basis $1,x,\dots,x^n$. Distinct polynomials of degree at most $n$ define distinct functions on $[0,1]$, because a nonzero polynomial has finitely many roots. Corollary~\ref{cor:norms} says that every finite-dimensional normed space is complete, so $(P_n,\|\cdot\|_\infty)$ is complete. Proposition~\ref{prop:complete}(a) says that a complete subset of a metric space is closed, so $P_n$ is closed in $C([0,1])$.

The interior is empty. Suppose some ball $B(f,r)$ were contained in $P_n$. Then $f\in P_n$, and $B(0,r)=B(f,r)-f$ would be contained in $P_n$. For an arbitrary $g\in C([0,1])$ with $g\ne f$, the function
\[h=f+\frac{r}{2\|g-f\|_\infty}(g-f)\]
satisfies $\|h-f\|_\infty=\frac r2$, so $h\in P_n$ and $h-f\in P_n$. Scaling puts $g-f$ in $P_n$, and then $g\in P_n$. The case $g=f$ is already inside $P_n$. This would mean $P_n=C([0,1])$.

That equality is false. The function $\gamma(x)=\bigl|x-\frac12\bigr|$ is continuous on $[0,1]$. Suppose a polynomial $p$ agreed with $\gamma$ on $[0,1]$. On $[\frac12,1]$ one has $\gamma(x)=x-\frac12$, so $p(x)-\bigl(x-\frac12\bigr)$ would vanish on an infinite set and $p(x)=x-\frac12$ as a polynomial. On $[0,\frac12]$ one has $\gamma(x)=\frac12-x$, so $x-\frac12=\frac12-x$ throughout that interval, hence $x=\frac12$. A single point is not the whole interval. So $\gamma\notin P_n$.

Each $P_n$ is closed and has empty interior, so each $P_n$ is nowhere dense. The set of all polynomial functions is $\bigcup_{n\ge0}P_n$, a countable union of nowhere dense sets, hence meagre. Corollary~\ref{cor:weierstrass} says that the same union is dense in $C([0,1])$. Theorem~\ref{thm:baire} puts a typical function outside every $P_n$. Polynomials are dense and meagre at the same time: they approximate every continuous function, and a typical continuous function is not one of them.
\end{solution}

\begin{solution}{ex:v7c05c}
\idea Vanishing on one fixed interval is a closed condition of empty interior, because adding or subtracting a small constant spoils identical vanishing. Rational endpoints reduce every subinterval to one of countably many cases.
\writeup Let $a<b$ be rational numbers in $[0,1]$, and let $Z_{a,b}$ be the set of $f\in C([0,1])$ with $f(x)=0$ for every $x\in[a,b]$. If $f_k\in Z_{a,b}$ and $f_k\to f$ uniformly, then $f=0$ on $[a,b]$, so $Z_{a,b}$ is closed.

The interior is empty. Fix $f$ and $\varepsilon>0$. The two functions $f+\frac\varepsilon2$ and $f-\frac\varepsilon2$ both lie at distance $\frac\varepsilon2$ from $f$. They cannot both vanish identically on $[a,b]$, because their difference is the constant $\varepsilon$, which is not identically zero. At least one of them lies outside $Z_{a,b}$ and inside the ball of radius $\varepsilon$ about $f$.

Suppose $f$ vanishes identically on some non-degenerate closed interval $[c,d]\subseteq[0,1]$. Choose rationals $a,b$ with $c<a<b<d$. Then $f$ vanishes on $[a,b]$, so $f\in Z_{a,b}$. The union of the sets $Z_{a,b}$ over all rational pairs $a<b$ in $[0,1]$ is therefore a countable union of nowhere dense sets. It is meagre. Theorem~\ref{thm:baire} says that a typical $f$ lies in none of the sets $Z_{a,b}$, and such an $f$ vanishes identically on no non-degenerate closed subinterval.
\end{solution}

\begin{solution}{ex:v7c05d}
\idea An even power kills every number strictly inside $(-1,1)$ and fixes the endpoints $\pm1$. The cosine equals $\pm1$ exactly on a discrete set of rationals, and those discrete sets grow, as $m$ grows, until they swallow every rational and no irrational.
\writeup (a) Fix $m$. If $|\cos(m!\pi x)|<1$, then $\cos(m!\pi x)^{2n}\to0$ as $n\to\infty$. If $|\cos(m!\pi x)|=1$, then the power equals $1$ for every $n$. The cosine of a real number has absolute value $1$ exactly when the argument is an integer multiple of $\pi$, so $m!\pi x=k\pi$ for some $k\in\Z$, that is, $x=\frac k{m!}$. Thus $g_m(x)=1$ on the set $\frac1{m!}\Z$ and $g_m(x)=0$ off that set. Each $f_{m,n}$ is continuous, so $g_m$ is a pointwise limit of continuous functions. Theorem~\ref{thm:baireclassone} is consistent with this: the continuity points of $g_m$ are the complement of $\frac1{m!}\Z$, which is open and dense, hence comeagre.

(b) Let $x=\frac ab$ be rational, with $b\ge1$. For every $m\ge b$ the factorial $m!$ is a multiple of $b$, so $m!\,x$ is an integer and $g_m(x)=1$. Therefore $\lim_m g_m(x)=1$. If $x$ is irrational, then $m!\,x$ is never an integer, because an integer value would make $x$ rational. Part~(a) gives $g_m(x)=0$ for every $m$, so the limit is $0$. The pointwise limit $\lim_m g_m$ is the indicator of $\Q$.

(c) At a rational the indicator equals $1$, and every neighbourhood contains irrationals, where the indicator equals $0$. At an irrational the indicator equals $0$, and every neighbourhood contains rationals, where it equals $1$. The indicator is discontinuous at every real number. Theorem~\ref{thm:baireclassone} says that a pointwise limit of continuous functions on the complete space $\R$ is continuous on a comeagre set, so it cannot be discontinuous everywhere. The indicator of $\Q$ is therefore not equal to any single pointwise limit of continuous functions. Part~(b) shows what two successive limits can do: each $g_m$ is one pointwise limit of continuous functions, and the indicator is the pointwise limit of the sequence $(g_m)$.
\end{solution}

\section*{Sources and further reading}
\addcontentsline{toc}{section}{Sources and further reading}

\begin{itemize}
\item For the split between null and meagre sets, read J.~C.~Oxtoby, \emph{Measure and Category}, 2nd ed., Springer, 1980, which treats Baire's theorem, the Banach--Mazurkiewicz theorem, the Liouville numbers, and the pairing of the two notions of small. It is the short account of Section~\ref{sec:categorymeasure}.
\item The thesis that introduces category, and proves the theorem on pointwise limits of continuous functions, is R.~Baire, \emph{Sur les fonctions de variables r\'eelles}, Annali di Matematica Pura ed Applicata 3 (1899), 1--123.
\item Theorem~\ref{thm:banachmaz} rests on two papers in the same journal: S.~Banach, \emph{\"Uber die Baire'sche Kategorie gewisser Funktionenmengen}, Studia Mathematica 3 (1931), 174--179, and S.~Mazurkiewicz, \emph{Sur les fonctions non d\'erivables}, same volume, 92--94.
\item The series for $T$ is T.~Takagi, \emph{A simple example of the continuous function without derivative}, Proceedings of the Physico-Mathematical Society of Japan 1 (1903), 176--177. Later identities are collected in P.~C.~Allaart and K.~Kawamura, \emph{The Takagi function: a survey}, Real Analysis Exchange 37 (2011/12), 1--54.
\item The quoted bound $ab\ge1$ is G.~H.~Hardy, \emph{Weierstrass's non-differentiable function}, Transactions of the American Mathematical Society 17 (1916), 301--325.
\item A short modern account of Baire's theorem, including the uses Volume~8 will take up, is W.~Rudin, \emph{Real and Complex Analysis}, 3rd ed., McGraw--Hill, 1987, Chapter~5.
\item Exercise~\ref{ex:polynomial} is a classical problem on flat derivatives. One place it is stated is W.~F.~Donoghue, \emph{Distributions and Fourier Transforms}, Academic Press, 1969, p.~53.
\item \emph{For the stories in the opening.} The lecture Weierstrass read on 18 July 1872 is printed in his \emph{Mathematische Werke}, vol.~2, Mayer \& M\"uller, 1895, 71--74. Hermite's letter of 20 May 1893, which we translated in the opening, is in \emph{Correspondance d'Hermite et de Stieltjes}, vol.~2, Gauthier-Villars, 1905. Poincar\'e's remarks on monsters and on the reasoning of our fathers are in \emph{Science et m\'ethode}, Flammarion, 1908, Book~II, Chapter~2. Khinchin's approximation theorem is A.~Khintchine, \emph{Einige S\"atze \"uber Kettenbr\"uche, mit Anwendungen auf die Theorie der Diophantischen Approximationen}, Mathematische Annalen 92 (1924), 115--125. The duality remark quotes W.~Sierpi\'nski, \emph{Sur la dualit\'e entre la premi\`ere cat\'egorie et la mesure nulle}, Fundamenta Mathematicae 22 (1934), 276--280, and P.~Erd\H{o}s, \emph{Some remarks on set theory}, Annals of Mathematics 44 (1943), 643--646.
\end{itemize}
